Un depósito contiene inicialmente $40\,\mathrm L$ de agua pura. En el modelo de mezcla se introduce una disolución con concentración suministrada de $3\,\mathrm{kg/L}$ a $2\,\mathrm{L/min}$ y sale mezcla homogénea a $3\,\mathrm{L/min}$. Suponga volúmenes aditivos y mezcla perfecta. Use $t=0$ al comenzar ambos flujos y limite el modelo al tiempo anterior al vaciado.

\begin{enumerate}[label=(\alph*),ref=\alph*,leftmargin=*]
\item Determine la cantidad de sal $S(t)$.
\item Calcule la cantidad cuando quedan $20\,\mathrm L$.
\item Determine cuándo es máxima la cantidad de sal.
\end{enumerate}
